
\documentclass{article}
\usepackage{colortbl}
\usepackage{makecell}
\usepackage{multirow}
\usepackage{supertabular}

\begin{document}

\newcounter{utterance}

\centering \large Interaction Transcript for game `cryptolect', experiment `basic', episode 0 with llama{-}3.1{-}8b{-}instruct.
\vspace{24pt}

{ \footnotesize  \setcounter{utterance}{1}
\setlength{\tabcolsep}{0pt}
\begin{supertabular}{c@{$\;$}|p{.15\linewidth}@{}p{.15\linewidth}p{.15\linewidth}p{.15\linewidth}p{.15\linewidth}p{.15\linewidth}}
    \# & \multicolumn{2}{c}{Player} && \multicolumn{2}{c}{Game Master} \\
    \hline

    \theutterance \stepcounter{utterance}  
    & & & \multicolumn{4}{p{0.6\linewidth}}{
        \cellcolor[rgb]{0.9,0.9,0.9}{
            \makecell[{{p{\linewidth}}}]{
                \texttt{\tiny{[P1$\langle$GM]}}
                \texttt{You are a field linguist. You have made contact with the only speaker of an unknown} \\
\texttt{language (the "informant"). Nobody has ever written down this language before. Your task} \\
\texttt{is to work out how it works and then translate sentences into it correctly.} \\
\\ 
\texttt{You will study the language by asking the informant to translate English sentences for} \\
\texttt{you. You may only talk about a small fixed world, using exactly these English words:} \\
\\ 
\texttt{Nouns: hunter, spear, stone, tree, hill, boat, river, lake} \\
\texttt{Adjectives: (this language has no adjectives)} \\
\texttt{Verbs (write them in this base form): carry, see, find, chase, feed} \\
\texttt{Articles: the (definite), a (indefinite)} \\
\texttt{Number words: two, three (no number word = one / singular)} \\
\texttt{These are ALL the words in play {-} the language has no other words.} \\
\\ 
\texttt{HOW TO WRITE AN ENGLISH SENTENCE} \\
\texttt{Every sentence is one clause with a subject, a verb, and an object, written in this fixed English notation:} \\
\texttt{    SUBJECT  verb  OBJECT} \\
\texttt{Each of SUBJECT and OBJECT is a noun phrase of the form:} \\
\texttt{    {[}the | a{]}  {[}number word{]}  {[}adjective ...{]}  noun} \\
\texttt{{-} 'the' = definite, 'a' = indefinite (singular only); leave the article off for an indefinite plural.} \\
\texttt{{-} a number word (two, three) sets the count; with no number word the noun is singular.} \\
\texttt{Examples of well{-}formed English sentences you may send:} \\
\texttt{    a hunter carry a spear} \\
\texttt{    three hunters carry a spear} \\
\\ 
\texttt{WHAT YOU ALREADY KNOW} \\
\texttt{The informant has already translated these sentences for you. Each line is} \\
\texttt{"<English>  =>  <the language>":} \\
\texttt{    a spear find three trees  =>  nelasu delide ki neso} \\
\texttt{    two stones carry three rivers  =>  tanisede doda terenode ki tana} \\
\texttt{    a stone see three stones  =>  tanise tanisede ki naka} \\
\texttt{    three trees chase two stones  =>  delide ki tanisede doda tedu} \\
\texttt{    a spear feed two rivers  =>  nelasu terenode doda loni} \\
\texttt{    two rivers see two hills  =>  terenode doda kalede doda naka} \\
\texttt{    a boat chase two rivers  =>  tunada terenode doda tedu} \\
\texttt{    three lakes find three spears  =>  lutide ki nelasude ki neso} \\
\texttt{    two hunters feed a spear  =>  sunode doda nelasu loni} \\
\\ 
\texttt{HOW THE SESSION WORKS} \\
\texttt{1. INVESTIGATION: You have 4 rounds of questioning. In each round you may ask} \\
\texttt{   the informant to translate up to 3 English sentences (so at most} \\
\texttt{   12 probes in total). The informant answers each one truthfully and} \\
\texttt{   exactly. Use these answers to figure out the words and the grammar (word order, how} \\
\texttt{   number/case/definiteness are marked, how words must agree, and any sound changes).} \\
\texttt{   Choose your probes deliberately {-} test specific hypotheses.} \\
\texttt{2. EXAM: After the investigation rounds, the informant will give you 10 NEW} \\
\texttt{   English sentences. You must translate each one INTO the language. Every answer is} \\
\texttt{   graded by EXACT match against how the informant would say it {-} every word, ending and} \\
\texttt{   space must be right (capitalisation and surrounding spaces are ignored). The exam} \\
\texttt{   sentences will use the same vocabulary but in combinations you have not been shown, so} \\
\texttt{   you must induce the general rules, not just memorise examples.} \\
\\ 
\texttt{RESPONSE FORMAT (for every investigation round)} \\
\texttt{List your probe sentences under a 'PROBES:' line, one per line. A THOUGHTS line is} \\
\texttt{optional {-} keep any reasoning brief so you always finish writing your PROBES:} \\
\texttt{THOUGHTS: <optional; brief>} \\
\texttt{PROBES:} \\
\texttt{1. <an English sentence, written as described above>} \\
\texttt{2. <another English sentence>} \\
\texttt{(up to 3 lines; at least one)} \\
\\ 
\texttt{Begin your first round of probes now.} \\
            }
        }
    }
     \\ \\

    \theutterance \stepcounter{utterance}  
    & \multicolumn{4}{p{0.6\linewidth}}{
        \cellcolor[rgb]{0.561,0.737,0.561}{
            \makecell[{{p{\linewidth}}}]{
                \texttt{\tiny{[P1$\rangle$GM]}}
                \texttt{PROBES:} \\
\texttt{1. a hunter see a spear} \\
\texttt{2. three hunters see a spear} \\
\texttt{3. a hunter see three spears} \\
\\ 
\texttt{I'm starting by testing the verb 'see' with different subjects and objects to see if there are any patterns in the word order or agreement.} \\
            }
        }
    }
    & & \\ \\

    \theutterance \stepcounter{utterance}  
    & & & \multicolumn{2}{p{0.3\linewidth}}{
        \cellcolor[rgb]{0.9,0.9,0.9}{
            \makecell[{{p{\linewidth}}}]{
                \texttt{\tiny{[GM$|$GM]}}
                \texttt{at most 3 probes per round, but 4 were given} \\
            }
        }
    }
    & & \\ \\

    \theutterance \stepcounter{utterance}  
    & & & \multicolumn{4}{p{0.6\linewidth}}{
        \cellcolor[rgb]{0.9,0.9,0.9}{
            \makecell[{{p{\linewidth}}}]{
                \texttt{\tiny{[P1$\langle$GM]}}
                \texttt{Your last response could not be processed: at most 3 probes per round, but 4 were given.} \\
\texttt{List your probes (a brief THOUGHTS line is optional):} \\
\texttt{PROBES:} \\
\texttt{1. <an English sentence>} \\
\texttt{(up to 3 probe lines; at least one)} \\
            }
        }
    }
     \\ \\

    \theutterance \stepcounter{utterance}  
    & \multicolumn{4}{p{0.6\linewidth}}{
        \cellcolor[rgb]{0.561,0.737,0.561}{
            \makecell[{{p{\linewidth}}}]{
                \texttt{\tiny{[P1$\rangle$GM]}}
                \texttt{THOUGHTS: Testing the verb 'see' with different subjects and objects.} \\
\\ 
\texttt{PROBES:} \\
\texttt{1. a hunter see a spear} \\
\texttt{2. three hunters see a spear} \\
\texttt{3. a hunter see three spears} \\
            }
        }
    }
    & & \\ \\

    \theutterance \stepcounter{utterance}  
    & & & \multicolumn{2}{p{0.3\linewidth}}{
        \cellcolor[rgb]{0.9,0.9,0.9}{
            \makecell[{{p{\linewidth}}}]{
                \texttt{\tiny{[GM$|$GM]}}
                {'kind': 'investigation', 'round': 0, 'probes_issued': 3, 'malformed': 0, 'probes': ['a hunter see a spear', 'three hunters see a spear', 'a hunter see three spears']}
            }
        }
    }
    & & \\ \\

    \theutterance \stepcounter{utterance}  
    & & & \multicolumn{4}{p{0.6\linewidth}}{
        \cellcolor[rgb]{0.9,0.9,0.9}{
            \makecell[{{p{\linewidth}}}]{
                \texttt{\tiny{[P1$\langle$GM]}}
                \texttt{TRANSLATIONS:} \\
\texttt{1. "a hunter see a spear"  =>  suno nelasu naka} \\
\texttt{2. "three hunters see a spear"  =>  sunode ki nelasu naka} \\
\texttt{3. "a hunter see three spears"  =>  suno nelasude ki naka} \\
\texttt{STATUS: investigation round 1 of 4 done. Probes used: 3 of 12. 3 investigation round(s) left before the exam (10 sentences).} \\
\texttt{Send your next probes now (a PROBES list; brief THOUGHTS optional).} \\
            }
        }
    }
     \\ \\

    \theutterance \stepcounter{utterance}  
    & \multicolumn{4}{p{0.6\linewidth}}{
        \cellcolor[rgb]{0.561,0.737,0.561}{
            \makecell[{{p{\linewidth}}}]{
                \texttt{\tiny{[P1$\rangle$GM]}}
                \texttt{THOUGHTS: The informant's translations suggest that the word order is subject{-}verb{-}object, and that the suffix "{-}de" is used to mark the object as plural.} \\
\\ 
\texttt{PROBES:} \\
\texttt{1. a hunter carry a stone} \\
\texttt{2. three hunters carry a stone} \\
\texttt{3. a hunter carry three stones} \\
            }
        }
    }
    & & \\ \\

    \theutterance \stepcounter{utterance}  
    & & & \multicolumn{2}{p{0.3\linewidth}}{
        \cellcolor[rgb]{0.9,0.9,0.9}{
            \makecell[{{p{\linewidth}}}]{
                \texttt{\tiny{[GM$|$GM]}}
                {'kind': 'investigation', 'round': 1, 'probes_issued': 3, 'malformed': 0, 'probes': ['a hunter carry a stone', 'three hunters carry a stone', 'a hunter carry three stones']}
            }
        }
    }
    & & \\ \\

    \theutterance \stepcounter{utterance}  
    & & & \multicolumn{4}{p{0.6\linewidth}}{
        \cellcolor[rgb]{0.9,0.9,0.9}{
            \makecell[{{p{\linewidth}}}]{
                \texttt{\tiny{[P1$\langle$GM]}}
                \texttt{TRANSLATIONS:} \\
\texttt{1. "a hunter carry a stone"  =>  suno tanise tana} \\
\texttt{2. "three hunters carry a stone"  =>  sunode ki tanise tana} \\
\texttt{3. "a hunter carry three stones"  =>  suno tanisede ki tana} \\
\texttt{STATUS: investigation round 2 of 4 done. Probes used: 6 of 12. 2 investigation round(s) left before the exam (10 sentences).} \\
\texttt{Send your next probes now (a PROBES list; brief THOUGHTS optional).} \\
            }
        }
    }
     \\ \\

    \theutterance \stepcounter{utterance}  
    & \multicolumn{4}{p{0.6\linewidth}}{
        \cellcolor[rgb]{0.561,0.737,0.561}{
            \makecell[{{p{\linewidth}}}]{
                \texttt{\tiny{[P1$\rangle$GM]}}
                \texttt{THOUGHTS: The informant's translations suggest that the suffix "{-}de" is used to mark the object as plural, and that the suffix "{-}ki" is used to mark the subject as plural.} \\
\\ 
\texttt{PROBES:} \\
\texttt{1. a hunter see a tree} \\
\texttt{2. three hunters see a tree} \\
\texttt{3. a hunter see the tree} \\
            }
        }
    }
    & & \\ \\

    \theutterance \stepcounter{utterance}  
    & & & \multicolumn{2}{p{0.3\linewidth}}{
        \cellcolor[rgb]{0.9,0.9,0.9}{
            \makecell[{{p{\linewidth}}}]{
                \texttt{\tiny{[GM$|$GM]}}
                {'kind': 'investigation', 'round': 2, 'probes_issued': 3, 'malformed': 0, 'probes': ['a hunter see a tree', 'three hunters see a tree', 'a hunter see the tree']}
            }
        }
    }
    & & \\ \\

    \theutterance \stepcounter{utterance}  
    & & & \multicolumn{4}{p{0.6\linewidth}}{
        \cellcolor[rgb]{0.9,0.9,0.9}{
            \makecell[{{p{\linewidth}}}]{
                \texttt{\tiny{[P1$\langle$GM]}}
                \texttt{TRANSLATIONS:} \\
\texttt{1. "a hunter see a tree"  =>  suno deli naka} \\
\texttt{2. "three hunters see a tree"  =>  sunode ki deli naka} \\
\texttt{3. "a hunter see the tree"  =>  suno deli naka} \\
\texttt{STATUS: investigation round 3 of 4 done. Probes used: 9 of 12. 1 investigation round(s) left before the exam (10 sentences).} \\
\texttt{Send your next probes now (a PROBES list; brief THOUGHTS optional).} \\
            }
        }
    }
     \\ \\

    \theutterance \stepcounter{utterance}  
    & \multicolumn{4}{p{0.6\linewidth}}{
        \cellcolor[rgb]{0.561,0.737,0.561}{
            \makecell[{{p{\linewidth}}}]{
                \texttt{\tiny{[P1$\rangle$GM]}}
                \texttt{THOUGHTS: The informant's translations suggest that the definite article "the" is not marked in the language, and that the word order is consistent with subject{-}verb{-}object.} \\
\\ 
\texttt{PROBES:} \\
\texttt{1. a hunter feed a spear} \\
\texttt{2. three hunters feed a spear} \\
\texttt{3. a hunter feed the river} \\
            }
        }
    }
    & & \\ \\

    \theutterance \stepcounter{utterance}  
    & & & \multicolumn{2}{p{0.3\linewidth}}{
        \cellcolor[rgb]{0.9,0.9,0.9}{
            \makecell[{{p{\linewidth}}}]{
                \texttt{\tiny{[GM$|$GM]}}
                {'kind': 'investigation', 'round': 3, 'probes_issued': 3, 'malformed': 0, 'probes': ['a hunter feed a spear', 'three hunters feed a spear', 'a hunter feed the river']}
            }
        }
    }
    & & \\ \\

    \theutterance \stepcounter{utterance}  
    & & & \multicolumn{4}{p{0.6\linewidth}}{
        \cellcolor[rgb]{0.9,0.9,0.9}{
            \makecell[{{p{\linewidth}}}]{
                \texttt{\tiny{[P1$\langle$GM]}}
                \texttt{TRANSLATIONS:} \\
\texttt{1. "a hunter feed a spear"  =>  suno nelasu loni} \\
\texttt{2. "three hunters feed a spear"  =>  sunode ki nelasu loni} \\
\texttt{3. "a hunter feed the river"  =>  suno tereno loni} \\
\\ 
\texttt{FINAL EXAM. Investigation is over. Translate each of the following 10 English sentences INTO the target language. Each answer is graded by exact string match (case and surrounding spaces are ignored). Every sentence uses only words you have already seen translated.} \\
\\ 
\texttt{1. three lakes feed a river} \\
\texttt{2. three stones carry two hunters} \\
\texttt{3. three spears chase a tree} \\
\texttt{4. a lake chase two boats} \\
\texttt{5. three trees see a boat} \\
\texttt{6. a hill see three hills} \\
\texttt{7. two rivers see a hunter} \\
\texttt{8. a tree see three hunters} \\
\texttt{9. a river see a boat} \\
\texttt{10. two hills see two trees} \\
\\ 
\texttt{Respond with the answers (you may add a brief THOUGHTS line first):} \\
\texttt{ANSWERS:} \\
\texttt{1. <translation of sentence 1>} \\
\texttt{... through 10.} \\
            }
        }
    }
     \\ \\

    \theutterance \stepcounter{utterance}  
    & \multicolumn{4}{p{0.6\linewidth}}{
        \cellcolor[rgb]{0.561,0.737,0.561}{
            \makecell[{{p{\linewidth}}}]{
                \texttt{\tiny{[P1$\rangle$GM]}}
                \texttt{THOUGHTS: Based on the informant's translations, I will use the patterns I've observed to translate the exam sentences.} \\
\\ 
\texttt{ANSWERS:} \\
\texttt{1. lutide ki tereno loni} \\
\texttt{2. tanisede ki sunode doda} \\
\texttt{3. nelasude ki deli tedu} \\
\texttt{4. lutide tereno doda tunada} \\
\texttt{5. deli ki tunada naka} \\
\texttt{6. kale naka kale naka} \\
\texttt{7. terenode doda suno naka} \\
\texttt{8. deli sunode ki naka} \\
\texttt{9. tereno tunada naka} \\
\texttt{10. kalede kalede ki neso} \\
            }
        }
    }
    & & \\ \\

    \theutterance \stepcounter{utterance}  
    & & & \multicolumn{2}{p{0.3\linewidth}}{
        \cellcolor[rgb]{0.9,0.9,0.9}{
            \makecell[{{p{\linewidth}}}]{
                \texttt{\tiny{[GM$|$GM]}}
                {'kind': 'exam', 'correct': 5, 'size': 10, 'results': [1, 0, 1, 0, 0, 0, 1, 1, 1, 0], 'token_similarity': 0.7611111111111112, 'answers': ['lutide ki tereno loni', 'tanisede ki sunode doda', 'nelasude ki deli tedu', 'lutide tereno doda tunada', 'deli ki tunada naka', 'kale naka kale naka', 'terenode doda suno naka', 'deli sunode ki naka', 'tereno tunada naka', 'kalede kalede ki neso'], 'targets': ['lutide ki tereno loni', 'tanisede ki sunode doda tana', 'nelasude ki deli tedu', 'luti tunadade doda tedu', 'delide ki tunada naka', 'kale kalede ki naka', 'terenode doda suno naka', 'deli sunode ki naka', 'tereno tunada naka', 'kalede doda delide doda naka']}
            }
        }
    }
    & & \\ \\

\end{supertabular}
}

\end{document}
